\subsubsection{Partícula en superficie de un cono}

Una partícula $P$ de masa $m$ se mueve sin roce sobre la superficie exterior de un cono de ángulo $\phi = \pi/4$. El sistema se encuentra en el espacio profundo (sin gravedad). La partícula comienza su movimiento a una distancia $r_0$ del vértice, con una velocidad inicial perpendicular al eje $\vec{e}_z$ y una velocidad angular $\dot{\theta}(t_0) = \omega_0$. Además de la fuerza normal, existe una fuerza de atracción central ejercida por el eje $\vec{e}_z$ sobre la partícula:
\[ \vec{f} = -\frac{B}{\rho^2} \vec{e}_\rho \]
donde $B$ es una constante conocida y $\rho$ es la coordenada radial cilíndrica. Considerando el origen en el vértice del cono, realice lo siguiente:

\begin{enumerate}
    \item Defina los marcos de referencia y sistemas de coordenadas (esféricas y cilíndricas) adecuados.
    \item Encuentre la velocidad angular $\dot{\theta}$ de $P$ en función de la coordenada esférica $r$.
    \item Determine si la fuerza $\vec{f}$ es o no conservativa.
    \item Escriba la energía mecánica total en términos de $\dot{r}$ y $r$.
    \item Analice si existen soluciones donde la coordenada $r$ esté acotada entre dos valores, $r_{min}$ y $r_{max}$.
\end{enumerate}

\begin{figure}[!ht]
    \centering
    \begin{tikzpicture}[>=Latex, scale=1.5]
        % --- Parámetros ---
        \def\Rbase{2}
        \def\Hcone{2}
        \def\ang{45} % pi/4

        % --- 1. Dibujo del Cono ---
        % Elipse de la base (atrás)
        \draw[dashed, gray] (\Rbase, 0) arc (0:180:{\Rbase} and 0.6);
        % Superficie lateral
        \draw[thick] (-\Rbase, 0) -- (0, \Hcone) -- (\Rbase, 0);
        % Elipse de la base (adelante)
        \draw[thick] (-\Rbase, 0) arc (180:360:{\Rbase} and 0.6);

        % --- 2. Ejes y Ángulo ---
        \draw[->, dashed] (0, \Hcone) -- (0, -0.5) node[below] {$\vec{e}_z$};
        \draw[dashed] (0, \Hcone) -- +(0.8, -0.8);
        \draw[thin] (0, \Hcone-0.4) arc (270:315:0.4);
        \node at (0.2, \Hcone-0.6) {$\frac{\pi}{4}$};

        % --- 3. Partícula y Vectores ---
        % Posición de la partícula P
        \coordinate (V) at (0, \Hcone); % Vértice
        \coordinate (P) at (1.2, {\Hcone - 1.2}); % Punto en la superficie
        
        \filldraw[shading=ball, ball color=black] (P) circle (2.5pt) node[right, xshift=3pt] {$P$};
        
        % Vector e_r (A lo largo del cono)
        \draw[->, thick] (P) -- +(135:-0.8) node[right] {$\vec{e}_r$};
        % Vector e_phi (Perpendicular a la superficie)
        \draw[->, thick] (P) -- +(45:0.8) node[right] {$\vec{e}_\phi$};
        
        % Líneas auxiliares para rho
        \draw[dashed, gray] (P) -- (0, {\Hcone - 1.2});
        \node[above] at (0.6, {\Hcone - 1.2}) {$\rho$};

    \end{tikzpicture}
    \caption{Partícula $P$ en movimiento sobre la superficie de un cono bajo una fuerza central atractiva.}
    \label{fig:cono_particula}
\end{figure}